\documentclass[11pt]{article}

% ============================================================
%  mu -> e gamma in warped (Randall--Sundrum) models
%  A slow, pedagogical internal review.
%  Written to track (a) the physics, (b) what is in OUR code,
%  (c) the subtleties and gaps.
%  This is an internal working note, not for distribution.
% ============================================================

\usepackage[margin=1in]{geometry}
\usepackage{amsmath,amssymb,amsthm}
\usepackage{mathtools}
\usepackage{graphicx}
\usepackage{booktabs}
\usepackage{enumitem}
\usepackage{xcolor}
\usepackage{framed}
\usepackage[colorlinks=true,linkcolor=blue,citecolor=blue,urlcolor=blue]{hyperref}

% lightweight "callout" box
\colorlet{shadecolor}{gray!12}
\newenvironment{callout}[1]{%
  \par\medskip\begin{shaded}\noindent\textbf{#1}\par\smallskip}{\end{shaded}\medskip}

\newcommand{\MKK}{M_{\mathrm{KK}}}
\newcommand{\LIR}{\Lambda_{\mathrm{IR}}}
\newcommand{\sw}{s_W}
\newcommand{\cw}{c_W}
\newcommand{\swsq}{s_W^2}
\newcommand{\cwsq}{c_W^2}
\newcommand{\half}{\tfrac{1}{2}}
\newcommand{\vev}{v}
\newcommand{\eps}{\varepsilon}
\newcommand{\YN}{Y_N}
\newcommand{\YNb}{\bar Y_N}
\newcommand{\YNmat}{\mathbf{Y}_N}
\newcommand{\YNbmat}{\bar{\mathbf{Y}}_N}
\newcommand{\spur}{(\YNb\YNb^\dagger)_{12}}
\newcommand{\BR}{\mathcal{B}}
\newcommand{\Vpmns}{V_{\mathrm{PMNS}}}
\newcommand{\code}[1]{\texttt{#1}}
\newcommand{\todo}[1]{{\color{red}[#1]}}

\theoremstyle{definition}
\newtheorem{definition}{Definition}

\title{\bf Lepton Flavour Violation: $\mu\to e\gamma$ in Warped Models\\[4pt]
\large A slow, pedagogical internal review --- physics, code, and subtleties}
\author{Internal working note (5D--Neutrino--Mixing repo)}
\date{\today}

\begin{document}
\maketitle

\begin{abstract}
\noindent
This note is a from-scratch, deliberately slow walk-through of the
charged-lepton-flavour-violating (cLFV) decay $\mu\to e\gamma$ and how it
constrains a Randall--Sundrum (RS) warped extra dimension. It is the single most
powerful lepton-sector probe in our flavour catalogue: a sharp, background-free
experimental limit (MEG~II, 2025) multiplied against a sharply model-dependent
RS prediction. The note assumes \emph{no} prior knowledge of dipole operators,
the seesaw mechanism, or ``RS-GIM'': all three are built up from the beginning.
Part~\ref{part:sm} defines the observable, the magnetic-dipole operator, and the
branching-ratio formula. Part~\ref{part:rs} explains \emph{why} warped models
generate $\mu\to e\gamma$, what the RS-GIM mechanism does, and where the
multi-TeV Kaluza--Klein (KK) floor comes from. Part~\ref{part:model} contrasts
the three model choices --- anarchic, Chen--Yu alignment, and the
lepton-minimal-flavour-violation (LMFV) hypothesis of Perez--Randall that our
code locks in --- and asks honestly whether our treatment is rigorous or proxy.
Part~\ref{part:code} documents exactly what our code does
(\code{flavorConstraints/muToEGamma.py}, the constraint \code{L001}, the
\code{LMFVLeptonParameters} carrier), including the \emph{critical honesty point}
that the veto is the branching ratio, not the famous coefficient $C=0.02$.
Part~\ref{part:subtle} catalogues the subtleties and the gap to ``rigorous.'' An
annotated reading guide closes the note.
\end{abstract}

\tableofcontents

% ====================================================================
\part*{Orientation: why this note exists}
\addcontentsline{toc}{section}{Orientation: why this note exists}
% ====================================================================

In our constraint catalogue, the quark-sector flavour-changing observables
($\eps_K$, $\Delta m_d$, $\Delta m_s$, $D^0$--$\bar D^0$) and $Z\to b\bar b$ are
the work-horses of the production scans. $\mu\to e\gamma$ is different in two
ways, and both are worth stating up front. First, it lives in the
\emph{lepton} sector: it is constraint \code{L001} in the \code{charged\_lepton}
family, it is \code{Severity.HARD}, and as of wave~W7 (commit \code{8ca55c8}) it
is \emph{live} in the catalogue --- but it is \emph{not} part of the quark-only
production scans, because those scans never build the lepton carrier it needs.
Second, its experimental side is essentially perfect (a single, clean
branching-ratio limit) while its theory side is among the most model-dependent
numbers in the whole catalogue: the prediction depends on the entire
right-handed-neutrino / seesaw sector (the neutrino Yukawa $\YN$, the PMNS
matrix, the Majorana scale $M_N$) through one off-diagonal spurion. To know what
our scan is really doing --- and what it would take to call \code{L001}
``rigorous'' --- we need to understand $\mu\to e\gamma$ from the ground up. That
is the purpose of this note.

A one-paragraph executive summary, to be unpacked over the following pages:

\begin{quote}
\small
$\mu\to e\gamma$ is a \emph{magnetic-dipole} flavour transition: a muon flips to
an electron while emitting a photon, which in the Standard Model (SM) is
forbidden to all practical purposes by the smallness of neutrino masses. Any new
physics that talks to both the muon and the electron generates an effective
dipole operator $\bar e\,\sigma^{\mu\nu}\mu\,F_{\mu\nu}$, and the decay rate
measures the square of its coefficient. In RS the photon, the leptons and the
heavy KK states all live in the warped bulk, so flavour-violating dipoles are
generated by KK exchange; their size is controlled by the \emph{differences} in
how the various lepton flavours localise (the ``RS-GIM'' mechanism). Whether the
prediction sits at, above, or below the experimental ceiling depends entirely on
the \emph{model of lepton flavour}. We adopt lepton minimal flavour violation
(LMFV, Perez--Randall): there, the only flavour-violating spurion is built from
the neutrino Yukawa, the relevant amplitude is $\propto\spur$, and naive
dimensional analysis (NDA) gives $\BR\sim 4\times10^{-8}\,|\spur|^2\,(3\,
\mathrm{TeV}/\MKK)^4$. Confronting this with the MEG~II limit
$\BR<1.5\times10^{-13}$ either pushes $\MKK$ into the multi-TeV range or forces
the spurion to be small --- which, because the spurion carries the PMNS matrix
and the seesaw scale, ties the LFV bound directly to the neutrino mass sector.
\end{quote}

\noindent The logical spine of the whole note, in one picture:
\begin{footnotesize}
\begin{verbatim}
  muon ---[ dipole operator  e mu sigma.F ]---> electron + photon   ( LFV: needs new physics )
                           |
   SM rate ~ (Dm_nu^2 / M_W^2)^2 ~ 10^-54   ==>  ANY signal = new physics, zero background
                           |
   RS: photon + leptons + Higgs in the bulk  ==>  KK exchange makes a dipole
                           |
   RS-GIM:  amplitude ~ DIFFERENCE of flavour localisations  ==>  small, but NOT zero
                           |
   size set by the MODEL of lepton flavour:
       |
       +-- anarchic        (ABP 0606021)  ==> M_KK floor ~ few-10 TeV  (strongest cLFV bound)
       |
       +-- Chen-Yu align   (degeneracy)   ==> dipole suppressed by alignment, floor lowered
       |
       +-- LMFV  (Perez-Randall 0805.4652) <== OUR CHOICE
                   spurion = (Ybar_N Ybar_N^dag)_12 ,  NDA coeff C ~ 0.02
                           |
   confront with MEG II 2025:  BR(mu->e gamma) < 1.5e-13   ==>  HARD veto in L001
                           |
   ( the model fork above, and a code-honesty caveat about C, come in Parts III-IV )
\end{verbatim}
\end{footnotesize}
\noindent Keep this map in mind; every box in it is unpacked below. The model
fork (which of the three flavour hypotheses our code locks in, and why) is
Part~\ref{part:model}; the most subtle implementation point --- that the coefficient
$C=0.02$ is \emph{not} what the code tests against --- is deferred to
\S\ref{sec:honesty}, where it can be proven from a unit test rather than merely
asserted.

\medskip
\noindent\textbf{How to read this note.} The audiences differ.
A reader here for the \emph{physics} should read Parts~\ref{part:sm}--\ref{part:model}
(the observable, the RS mechanism, the model fork) and may skim Part~\ref{part:code}
(the code documentation) on a first pass, returning to it only for the honesty
point \S\ref{sec:honesty}. A reader here for the \emph{implementation} can move
quickly through Parts~\ref{part:sm}--\ref{part:rs} and dwell on
Parts~\ref{part:code}--\ref{part:subtle}. Everyone should read \S\ref{sec:honesty}
and the worked example in \S\ref{sec:floor}.

% ====================================================================
\part{The Standard-Model framework: what $\mu\to e\gamma$ is}
\label{part:sm}
% ====================================================================

\section{The observable, in words and in numbers}
\label{sec:observable}

A free positive muon decays overwhelmingly through $\mu^+\to e^+\nu_e\bar\nu_\mu$
(Michel decay). The process of interest here is the rare \emph{radiative}
variant in which the two neutrinos are absent:
\begin{equation}
\mu^+ \to e^+ \gamma .
\end{equation}
The two final-state neutrinos are what carry away lepton-family number in the
ordinary decay; removing them means the muon-ness of the initial state has been
converted directly into electron-ness. That is \emph{charged lepton flavour
violation} (cLFV): a transition between lepton generations with no compensating
neutrinos. The experimental signature is clean and almost background-free: a
back-to-back, monochromatic $e^+$ and $\gamma$, each carrying half the muon rest
energy, $E_e = E_\gamma = m_\mu/2 \approx 52.8$~MeV.

\begin{definition}[Branching ratio]
The branching ratio $\BR(\mu\to e\gamma)$ is the fraction of all muon decays that
proceed through the $e\gamma$ final state,
\[
\BR(\mu\to e\gamma) \equiv \frac{\Gamma(\mu\to e\gamma)}{\Gamma(\mu\to\text{all})}
   \simeq \frac{\Gamma(\mu\to e\gamma)}{\Gamma(\mu\to e\nu\bar\nu)},
\]
the last step because Michel decay saturates the total muon width. It is the
quantity experiments report as an upper limit.
\end{definition}

\section{Why the SM rate is, for our purposes, exactly zero}
\label{sec:smrate}

In the SM augmented with massive neutrinos, $\mu\to e\gamma$ \emph{is} generated:
at one loop, a $W$ boson and a neutrino circulate, and neutrino oscillations
($\nu_\mu \leftrightarrow \nu_e$ mixing) supply the flavour change. But the
amplitude is throttled by the GIM mechanism: it is proportional to the
\emph{difference} of neutrino-mass-squared splittings over the $W$ mass squared,
\begin{equation}
\BR(\mu\to e\gamma)_{\rm SM} \sim \frac{3\alpha}{32\pi}
   \left|\sum_i U^*_{ei}U_{\mu i}\frac{\Delta m^2_{i1}}{M_W^2}\right|^2
   \;\sim\; 10^{-54}.
\label{eq:smrate}
\end{equation}
The number is so far below any conceivable experiment (the MEG~II reach
$\sim 1.5\times10^{-13}$ is some $41$ orders of magnitude --- a factor
$\sim 10^{41}$ --- above this SM rate) that for catalogue purposes the SM
prediction is \emph{zero}. This is the single most important practical feature of
$\mu\to e\gamma$: there is no irreducible SM background to subtract, no hadronic
uncertainty, no theory error on the ``floor.'' \textbf{Any} observed signal is
unambiguously new physics. Our code encodes exactly this: \code{L001} sets
\code{sm\_prediction = 0.0} and applies its entire budget to the
\emph{new-physics} branching fraction (see \S\ref{sec:codeL001}).

\section{The effective operator: a magnetic dipole}
\label{sec:operator}

Because the photon is on-shell and massless, the most general
flavour-violating $\mu$--$e$--$\gamma$ vertex is fixed by gauge invariance to be
a \emph{magnetic-dipole} (chirality-flipping) operator. After integrating out
whatever heavy physics generates it, the relevant piece of the Lagrangian is
\begin{equation}
\mathcal{L}_{\rm dipole} \;=\;
   \frac{e\,m_\mu}{2}\;\bar e\,\sigma^{\mu\nu}\big(A_L P_L + A_R P_R\big)\mu\;
   F_{\mu\nu} \;+\; \text{h.c.},
\label{eq:dipoleop}
\end{equation}
where $\sigma^{\mu\nu}=\tfrac{i}{2}[\gamma^\mu,\gamma^\nu]$,
$P_{L,R}=\tfrac{1}{2}(1\mp\gamma_5)$ are chirality projectors, $F_{\mu\nu}$ is the
photon field strength, and $A_{L,R}$ are dimensionful Wilson coefficients
(dimension $1/\text{mass}^2$) carrying the model dependence. Three features of
\eqref{eq:dipoleop} are worth pausing on:
\begin{itemize}[leftmargin=1.4em]
\item \textbf{It flips chirality.} $\bar e\,\sigma^{\mu\nu}\mu$ connects a
left-handed lepton to a right-handed one, which is why the explicit factor of
$m_\mu$ (a chirality-flipping mass insertion) appears out front. A theory with no
source of chirality breaking cannot generate it.
\item \textbf{It is a dimension-6 operator} (after counting the Higgs vev that
supplies $m_\mu$): suppressed by two powers of the new-physics scale. Hence the
amplitude scales as $1/\MKK^2$ and the \emph{rate} as $1/\MKK^4$. This is the
origin of the famous $(\MKK)^{-4}$ in every branching-ratio formula below.
\item \textbf{It is loop-induced.} A dipole requires a photon to attach to a
charged line inside a loop; there is no tree-level dipole. The loop supplies a
factor $1/16\pi^2$, which the NDA estimate will track.
\end{itemize}
The decay rate is the square of \eqref{eq:dipoleop}'s coefficient; in the
standard normalisation,
\begin{equation}
\BR(\mu\to e\gamma) = \frac{48\pi^3\alpha}{G_F^2}
   \left(|A_L|^2 + |A_R|^2\right),
\label{eq:brgeneric}
\end{equation}
so the entire job of any model --- SM, supersymmetry, or RS --- is to compute the
two numbers $A_L, A_R$.

\section{The experimental constraint: MEG II 2025}
\label{sec:expt}

The strongest existing limit comes from the MEG~II experiment at PSI, which
stops $\sim 10^{14}$ positive muons and looks for the back-to-back $e^+\gamma$
signature. The value our catalogue uses (the primary anchor in \code{L001.yaml})
is
\begin{equation}
\boxed{\;\BR(\mu^+\to e^+\gamma) < 1.5\times10^{-13}\quad (90\%\ \text{C.L.}),\;}
\label{eq:megii}
\end{equation}
from the 2021--2022 physics runs (published October 2025; quoted search
sensitivity $2.2\times10^{-13}$). For context, the catalogue sidecar also stores
the experimental history, which is a thirty-year march of improving limits:

\begin{center}
\begin{tabular}{lll}
\toprule
Experiment / dataset & $\BR$ upper limit (90\% C.L.) & era \\
\midrule
MEGA (used by Perez--Randall) & $1.2\times10^{-11}$ & paper-era input \\
MEG full 2009--2013          & $4.2\times10^{-13}$ & 2016 \\
MEG~II 2021-only             & $7.5\times10^{-13}$ & 2024 \\
MEG + MEG~II combination     & $3.1\times10^{-13}$ & 2024 \\
\textbf{MEG~II 2021--2022 (primary)} & $\mathbf{1.5\times10^{-13}}$ & \textbf{2025} \\
\bottomrule
\end{tabular}
\end{center}

\noindent The first row matters for our code: when Perez--Randall did their NDA
estimate in 2008 they used the contemporaneous MEGA bound
$1.2\times10^{-11}$, and \emph{that} input is where the coefficient $0.02$ comes
from (\S\ref{sec:Cnumber}). The last row is what the live scan tests against. The
factor of $\sim 80$ between them ($1.2\times10^{-11}$ vs $1.5\times10^{-13}$)
tightens any derived coefficient by $\sqrt{80}\approx 9$ and is the reason the
``paper'' coefficient and the ``live'' coefficient differ.

% ====================================================================
\part{How warped models generate $\mu\to e\gamma$}
\label{part:rs}
% ====================================================================

\section{RS recap: where leptons and the photon live}
\label{sec:rsrecap}

Recall the geometry (repo conventions). Spacetime is a slice of AdS$_5$ between a
UV brane at $z=1/k$ (Planck scale) and an IR brane at $z=1/\LIR$ (TeV scale); the
warp factor is $\eps=\LIR/k\sim10^{-15}$. A 5D fermion carries a dimensionless
bulk-mass parameter $c=M_5/k$ that fixes where its massless zero mode lives:
\begin{itemize}[leftmargin=1.4em]
\item $c>\half$: the zero mode is peaked at the \emph{UV} brane, with an
exponentially small overlap at the IR brane;
\item $c<\half$: the zero mode is peaked at the \emph{IR} brane, with an
$\mathcal{O}(1)$ overlap there.
\end{itemize}
The normalised brane overlap (an ``$f$-factor'') is, in the repo's convention,
$f_{\rm IR}^2(c)=(\half-c)/(1-\eps^{1-2c})$. Light leptons, which need small
masses, are UV-localised ($c>\half$, small $f$); this is the central RS idea ---
hierarchies come from \emph{geometry}, not from hierarchical 5D couplings.

The photon (and the $W$, $Z$, gluon) are bulk gauge fields. Their zero modes are
the SM gauge bosons; their excited KK modes are localised toward the IR brane
with masses $m_n \simeq x_n\LIR$, where the $x_n$ are Bessel roots
(the first gauge KK root for Neumann--Neumann boundary conditions is
$x_1\approx 2.45$, the constant \code{GAUGE\_KK\_ROOT\_NN = 2.448687135269161}
in our code). \textbf{This factor of $2.45$ matters for every TeV bound we quote}
and is a recurring source of confusion --- see \S\ref{sec:Mconvention}.

\section{The mechanism: KK exchange and the chirality flip}
\label{sec:mechanism}

How does a dipole appear? The Higgs sits on (or near) the IR brane, and the heavy
KK fermions and KK gauge bosons pile up there too. A light external muon
(UV-localised) has a small but non-zero tail reaching the IR brane; there it can
(i) pick up a chirality-flipping mass insertion from the Higgs vev, (ii) emit a
photon, and (iii) be converted to an electron by a flavour-violating
KK-fermion/KK-gauge exchange. Diagrammatically this is exactly the loop that
generates \eqref{eq:dipoleop}: a heavy KK line inside, a photon attached, a mass
insertion to flip chirality. Integrating out the KK tower gives a dipole
coefficient
\begin{equation}
A \;\sim\; \frac{1}{16\pi^2}\,\frac{1}{\MKK^2}\times(\text{flavour structure}),
\label{eq:Aparam}
\end{equation}
the $1/16\pi^2$ from the loop and the $1/\MKK^2$ from the heavy propagators ---
the universal decoupling. As Agashe, Blechman, and Petriello
(hep-ph/0606021) stress, the calculation is only \emph{finite and trustworthy}
when the Higgs propagates in the bulk; with a strictly brane-localised Higgs the
dipole is formally ultraviolet-sensitive (uncalculable), and one must treat its
coefficient as an NDA estimate. Our code does the latter --- it uses an NDA
coefficient, not a finite loop integral --- which is the honest reason
\code{L001}'s \emph{prediction} side is model-dependent
(\S\ref{sec:rigorous}).

\section{RS-GIM: why the answer is small but not zero}
\label{sec:rsgim}

The ``flavour structure'' in \eqref{eq:Aparam} is where the magic --- and the
constraint --- lives. In RS, the flavour-violating coupling of a KK gauge boson to
two leptons is governed by their bulk localisations. If \emph{all} lepton flavours
sat at exactly the same place in the bulk, the KK gauge couplings would be
flavour-universal and rotate away: no flavour violation, in perfect analogy with
the ordinary GIM mechanism --- it is the \emph{same} idea as the
mass-splitting-\emph{difference} structure that suppresses the SM rate
\eqref{eq:smrate}, where the amplitude vanishes in the limit of degenerate
neutrino masses. They do \emph{not} sit at the same place --- that is
how the mass hierarchy is generated --- but they sit at \emph{similar} places (all
light leptons are UV-localised), so the flavour-violating piece is the small
\emph{difference} of their localisations. The reason a \emph{difference} is the
violation is basis misalignment: the KK gauge couplings are diagonal in the
\emph{flavour-localisation} basis (each generation couples according to where its
zero mode sits in the bulk), but the physical states live in the \emph{mass}
basis; the mismatch between these two bases --- set by how differently the
generations are localised --- is what cannot be rotated away. This is the
\textbf{RS-GIM mechanism}:
\begin{equation}
A_{\mu\to e} \;\propto\; (\alpha_\mu - \alpha_e),
\qquad \alpha_i = \text{(KK coupling set by the localisation of flavour } i),
\end{equation}
where the $\alpha_i$ are controlled by the same $f$-factors $f_{\rm IR}(c)$ of
\S\ref{sec:rsrecap} that set the lepton masses. The individual $\alpha_i$ are all
$\mathcal{O}(0.1)$ in the anarchic estimate of hep-ph/0606021; the suppression is
\emph{not} in their magnitudes but in their \emph{difference}
$(\alpha_\mu-\alpha_e)$, which is what enters the dipole. (This
$(\alpha_\mu-\alpha_e)$ picture is the \emph{anarchic} flavour-violation
mechanism of \S\ref{sec:anarchic}; it is distinct from the single LMFV spurion
$\spur$ our code actually uses, built from neutrino Yukawas rather than from the
charged-lepton localisations $\alpha_i$ --- the two are different sources of
$\mu\to e$ violation, and the LMFV one is what \eqref{eq:prBR} encodes.) The same
geometry that explains \emph{why} the electron is lighter than the muon also
\emph{suppresses} (but does not forbid) $\mu\to e$ transitions. RS-GIM is the
lepton-sector analogue of the quark-sector RS-GIM that protects $\eps_K$ and
$\Delta m_K$. It is what makes the model survive at all; whether it survives
comfortably or barely is the quantitative question of Part~\ref{part:model}.

\section{The branching-ratio formula our code uses (LMFV/NDA)}
\label{sec:brformula}

Specialising to the LMFV hypothesis (motivated in \S\ref{sec:lmfv}), the only
flavour-violating object is the neutrino Yukawa, and the $\mu$--$e$ dipole is
controlled by the (1,2) entry of the Hermitian spurion $\YNb\YNb^\dagger$. The
Perez--Randall NDA dipole operator (their Eq.~(29)) and the resulting
branching ratio (their Eq.~(30), with the spurion bound their Eq.~(31)) are
\begin{equation}
\mathcal{O}^{\rm IR}_{\rm dipole} \approx
   \frac{e\,m_\mu}{2\cdot 16\pi^2\,k^2}\,\spur\,
   F^{\mu\nu}\,\bar e_L\,\sigma_{\mu\nu}\,\mu,
\label{eq:prdipole}
\end{equation}
\begin{equation}
\boxed{\;\BR(\mu\to e\gamma) \simeq 4\times10^{-8}\;
   |\spur|^2\;\left(\frac{3\ \mathrm{TeV}}{\MKK}\right)^{4}.\;}
\label{eq:prBR}
\end{equation}
Here $\YNb = 2k\,\YN$ is the dimensionless rescaled neutrino Yukawa (the same
$2k$ rescaling the repo uses everywhere), so $\spur$ is dimensionless and
$\mathcal{O}(1)$ for anarchic Yukawas. Two structural facts to carry forward:
\begin{itemize}[leftmargin=1.4em]
\item The prefactor $4\times10^{-8}$ is the NDA packaging of
$(\alpha/\ldots)(1/16\pi^2)^2$ etc., evaluated at the reference scale $3$~TeV. In
our code it is the constant \code{PREFAC\_BR = 4.0e-8}.
\item The $(\MKK)^{-4}$ scaling is the dipole-rate scaling of
\S\ref{sec:operator}: raising the KK scale by a factor of two suppresses the rate
by sixteen. This is what converts an experimental limit into a \emph{lower bound}
on $\MKK$.
\end{itemize}

\section{Where the $\MKK$ floor comes from}
\label{sec:floor}

Turn \eqref{eq:prBR} around. Imposing $\BR<\BR_{\rm lim}$ gives, equivalently, a
bound on the spurion or a bound on the KK scale:
\begin{equation}
|\spur| \;\lesssim\; \sqrt{\frac{\BR_{\rm lim}}{4\times10^{-8}}}
   \left(\frac{\MKK}{3\ \mathrm{TeV}}\right)^{2}
   \;\equiv\; C\left(\frac{\MKK}{3\ \mathrm{TeV}}\right)^{2},
\label{eq:spurbound}
\end{equation}
which \emph{defines} the coefficient $C$ (this is the form
\code{coefficient\_from\_br\_limit} in our code computes). For an
$\mathcal{O}(1)$ anarchic spurion $|\spur|\sim 1$, \eqref{eq:spurbound} forces
$\MKK \gtrsim 3\,\mathrm{TeV}/\sqrt{C}$. With the paper-era input
$C\approx 0.02$ this is $\MKK\gtrsim 3/\sqrt{0.02}\approx 21$~TeV; with the
MEG~II input the coefficient drops to $C\approx 1.9\times10^{-3}$ and the floor
climbs further still. This is the celebrated \textbf{RS lepton-LFV problem}: in
the anarchic case the lepton-flavour-violating observables collectively force
$\MKK$ to many TeV, well above the electroweak floor set by $S,T$. ABP make this
case from the \emph{full set} of cLFV observables --- $\mu$--$e$ conversion in
nuclei, rare $\mu$ decays (including $\mu\to e\gamma$), and rare $\tau$ decays ---
and emphasise a built-in \emph{tension} between the loop-induced dipole
($\mu\to e\gamma$) and tree-level $\mu$--$e$ conversion, which depend oppositely on
the Yukawa size, so that no single observable can be tuned away without aggravating
another. (In much of the RS-LFV literature $\mu$--$e$ conversion, not
$\mu\to e\gamma$, is the leading bound.) The LMFV and alignment models of
Part~\ref{part:model} exist precisely to relax this collective tension.

\begin{callout}{Worked numerical example --- one benchmark point, all the numbers.}
This is the actual benchmark pinned in our unit test
(\code{tests/.../test\_L001.py}). (The spurion $\spur$ used here is fully defined
only in \S\ref{sec:lmfv}--\S\ref{sec:brformula}: it is the $(1,2)$ entry of the
Hermitian matrix $\YNb\YNb^\dagger$ with $\YNb$ the PMNS-rotated rescaled neutrino
Yukawa. For this example treat $|\spur|$ as a given number.) The point has spurion
magnitude $|\spur|\approx 0.0623$ at $\MKK=3$~TeV (so the reference factor
$(3\,\mathrm{TeV}/\MKK)^4=1$); this is the test's
$\code{\_ORACLE\_LHS}=0.062266\ldots$, the absolute value of
$\code{\_ORACLE\_OFFDIAG}=-0.0348+0.0517\,i$. Note this is already an
order of magnitude \emph{below} $\mathcal{O}(1)$ --- the LMFV partial
alignment and the small neutrino Yukawa eigenvalues have done real work here, and
this is \emph{not} the raw anarchic $|\spur|\sim 1$ scenario of
\S\ref{sec:anarchic}. Even so, the point fails. The predicted new-physics
branching ratio \eqref{eq:prBR} is
\[
\BR_{\rm NP} = 4\times10^{-8}\times(0.0623)^2\times 1
   \;\approx\; 1.55\times10^{-10}
   \quad(\text{the test's }\code{\_ORACLE\_BR}=1.5508\ldots\times10^{-10}).
\]
Confront with MEG~II, $\BR_{\rm lim}=1.5\times10^{-13}$:
\[
\frac{\BR_{\rm NP}}{\BR_{\rm lim}} = \frac{1.55\times10^{-10}}{1.5\times10^{-13}}
   \;\approx\; 1.03\times10^{3}
   \quad(\text{the test's }\code{\_ORACLE\_BR\_RATIO}=1033.9),
\]
so this point \emph{fails} the bound by three orders of magnitude --- it is
\textbf{vetoed}. (For comparison the inert dipole diagnostic at the paper
coefficient is $\text{rhs}=C(\MKK/3000)^2=0.02$ and the dipole ratio
$|\spur|/\text{rhs}=3.11$, the test's
$\code{\_ORACLE\_DIP\_RATIO\_PAPER\_C}=3.1133$ --- but recall \S\ref{sec:honesty},
that ratio is \emph{not} the gate.) To rescue the point one would need to raise
$\MKK$ by $(1.03\times10^3)^{1/4}\approx 5.7$, i.e.\ to $\MKK\approx 17$~TeV, or to
suppress the spurion further. The same arithmetic, in words: \emph{even a
sub-$\mathcal{O}(1)$ spurion ($|\spur|\approx 0.06$) at $3$~TeV overshoots the
MEG~II ceiling by $\sim 1000\times$, and only a $\sim 6\times$ heavier KK scale
(or a $\sim 30\times$ smaller spurion) brings it under.}
\end{callout}

% ====================================================================
\part{Model dependence: anarchic vs.\ Chen--Yu vs.\ LMFV}
\label{part:model}
% ====================================================================

\emph{The experimental limit \eqref{eq:megii} is model-independent. The
\emph{prediction} \eqref{eq:prBR} is not: the spurion $\spur$ and the relevant
$C$ depend entirely on the assumed model of lepton flavour. This part lays out
the three options and explains which one our code locks in, and why.}

\section{Option 1: anarchic flavour (the maximal, most constrained case)}
\label{sec:anarchic}

The default RS flavour hypothesis is \emph{anarchy}: the 5D Yukawa couplings are
structureless $\mathcal{O}(1)$ complex numbers, and \emph{all} the observed
hierarchy comes from the $f$-factors. This is the framework of
Agashe--Blechman--Petriello (hep-ph/0606021), whose central thesis is that
precisely because flavour is geometric, LFV becomes a \emph{direct probe of the
geometry}: the same localisation pattern that fits the charged-lepton masses
predicts the $\mu\to e$ dipole, with no freedom left to hide it. In the anarchic
case the spurion is generically $\mathcal{O}(1)$, the individual $\alpha_i$ are
only $\mathcal{O}(0.1)$ (RS-GIM acts through their \emph{differences}, not their
magnitudes), and the bound \eqref{eq:spurbound} bites hard: the cLFV observables
collectively push the lightest KK gauge mass into the few-to-ten-TeV range across
the natural parameter space. ABP frame this as a \emph{virtue}: forthcoming
$\mu\to e\gamma$ and $\mu$-conversion experiments ``definitively test'' the RS
geometric origin of the lepton hierarchy. The price is that the minimal anarchic
model is in tension with EW precision unless $\MKK$ is uncomfortably high.

\section{Option 2: generic alignment / degeneracy (the ``Chen--Yu'' option)}
\label{sec:chenyu}

A second route --- which we label generically, as it is a mechanism rather than a
single canonical reference --- is to give up anarchy and \emph{align} (or partially
degenerate) the lepton bulk masses so that the flavour-violating differences
$(\alpha_\mu-\alpha_e)$ in the RS-GIM expression are suppressed by construction.
If the relevant lepton localisations are forced to be degenerate, the dipole
spurion shrinks and the $\mu\to e\gamma$ bound is relaxed --- much as
\emph{quark}-sector bulk-mass alignment relaxes $\eps_K$. The cost is that the
alignment is an extra assumption (a flavour symmetry or a boundary condition) not
required by the geometry itself; it trades predictivity for a lower KK floor. Our
internal model survey recorded this as a viable but less clean alternative: it
works, but it requires importing structure that anarchy and LMFV avoid.

\section{Option 3: LMFV (Perez--Randall) --- our locked choice}
\label{sec:lmfv}

The model our code adopts is \emph{lepton minimal flavour violation} (LMFV),
following Perez--Randall (0805.4652). The idea is the lepton analogue of minimal
flavour violation in the quark sector: \emph{assume the only spurions that break
lepton flavour are the Yukawa matrices themselves}. Concretely, Perez--Randall
take the bulk mass matrices (the $c$-parameters promoted to matrices) to be
functions of the Yukawas,
\begin{align}
C_E &= \mathbb{1} + a\,k^2\,Y_E^\dagger Y_E, \nonumber\\
C_N &= \mathbb{1} + b\,k^2\,Y_N^\dagger Y_N, \\
C_L &= \mathbb{1} + a_E\,k^2\,Y_E Y_E^\dagger + b_N\,k^2\,Y_N Y_N^\dagger,
\nonumber
\end{align}
with $a,b,a_E,b_N$ real $\mathcal{O}(1)$ coefficients. Physically, building the
$c$-parameter matrices out of the Yukawas forces them to be (approximately)
diagonalised by the \emph{same} unitary rotations as the Yukawas themselves; the
localisation basis and the mass basis are then nearly aligned by construction, so
the leftover misalignment that feeds the dipole is small and calculable rather than
anarchic. The consequence is \emph{partial alignment}: because the bulk masses are
built from the same Yukawas that produce the fermion masses, the flavour-violating
misalignment is not arbitrary but is governed by the Yukawa spurions. For the charged-lepton dipole,
working in the charged-lepton mass basis, the only surviving flavour-violating
combination relevant to $\mu\to e\gamma$ is
\begin{equation}
\spur, \qquad \YN \to \Vpmns\cdot\mathrm{diag}(Y_{N_1},Y_{N_2},Y_{N_3}),
\label{eq:YNrot}
\end{equation}
i.e.\ the neutrino Yukawa rotated into the physical basis by the PMNS matrix.
This is the precise statement that \textbf{$\mu\to e\gamma$ in LMFV is a window on
the right-handed-neutrino / seesaw sector}: the spurion carries the PMNS angles
(through $\Vpmns$), the neutrino Yukawa eigenvalues $Y_{N_i}$, and --- because
those eigenvalues are fixed by the seesaw relation
$m_{\nu_i}\propto Y_{N_i}^2/M_N$ --- the Majorana scale $M_N$. Change the neutrino
spectrum, ordering, or Majorana phases, and the $\mu\to e\gamma$ prediction moves.

Why LMFV and not the other two? Our internal model survey (recorded in the
orchestration ledger) compared LMFV against anarchic, Chen--Yu, and a gauged
variant and concluded LMFV is the \emph{cleanest}: it is predictive (no extra
flavour symmetry imposed by hand, unlike Chen--Yu), it is less constrained than
raw anarchy (the partial alignment softens the dipole), and its NDA coefficient
is already coded. The verdict recorded was: ``model locked $=$ LMFV
(Perez--Randall), spurion $\spur$, $C=0.02$ NDA; NDA conservative \emph{within}
LMFV; bounds rescale $\sim 3\times$ at MEG~II.''

\begin{callout}{Common confusions, defused.}
\begin{itemize}[leftmargin=1.2em,nosep]
\item[\textbf{A.}] \emph{``The coefficient $C=0.02$ is what the code tests
against.''} No --- the gate is the branching ratio $\BR_{\rm NP}\le\BR_{\rm lim}$.
$C$ is an \emph{inert diagnostic} that scales only the dipole right-hand side, not
the pass/fail. A unit test pins this (\S\ref{sec:honesty}).
\item[\textbf{B.}] \emph{``$\mu\to e\gamma$ is a quark-sector constraint in our
scans.''} No --- it is \code{L001} in the \code{charged\_lepton} family, and it is
\emph{not} evaluated in the quark-only production scans, which never build the
lepton carrier. It returns ``NOT EVALUATED'' there (\S\ref{sec:notinscan}).
\item[\textbf{C.}] \emph{``The bound depends only on the lepton masses.''} No ---
in LMFV it depends on the entire neutrino sector: $\YN$, the PMNS matrix, the
Majorana scale $M_N$, the ordering and the Majorana phases, all through
$\spur$ \eqref{eq:YNrot}.
\item[\textbf{D.}] \emph{``$\MKK$ means one fixed thing.''} It is overloaded by
the factor $x_1\approx 2.45$: the physical first KK gauge mass is
$x_1\LIR$, but the repo's default LFV convention uses $\MKK=\LIR$
(\S\ref{sec:Mconvention}).
\item[\textbf{E.}] \emph{``The SM contributes a background.''} No --- the SM rate
is $\sim 10^{-54}$, so \code{L001} sets \code{sm\_prediction = 0.0} and the bound
is a pure new-physics veto.
\end{itemize}
\end{callout}

\section{Does custodial symmetry or alignment change \emph{this} bound?}
\label{sec:custodial}

A natural question, given the importance of custodial symmetry for the oblique
$S,T$ bounds: does it help with $\mu\to e\gamma$? The honest answer is
\emph{essentially no, not directly}. Custodial $SU(2)_R\times P_{LR}$ is designed
to protect the \emph{electroweak} observables (the $T$ parameter, the
$Z\bar b_L b_L$ vertex); it says nothing about the lepton-flavour-violating dipole
spurion. The custodial route to a $\sim 3$~TeV electroweak KK floor --- extending
the bulk gauge group to a left--right symmetric $SU(2)_L\times SU(2)_R$ --- is a
known solution that ABP note (attributing it to earlier work) but do \emph{not}
themselves scan; it addresses the oblique and $Z\bar b_L b_L$ bounds, not the
dipole. The physics point stands on its own logic, not on a quoted ABP scan:
custodial symmetry can let the KK scale be light \emph{as far as $S,T$ are
concerned}, but the LFV dipole spurion is untouched, so $\mu\to e\gamma$ (and
$\mu$--$e$ conversion) then become the binding constraints instead. The lever that
actually moves the LFV bound is the
\emph{flavour} structure --- anarchy vs.\ alignment vs.\ LMFV --- not the
electroweak custodial structure. This is why our LFV treatment is a separate
constraint (\code{L001}) with its own model choice (LMFV), decoupled from the
custodial \code{ew\_model} switch that governs the oblique and $Zb\bar b$
constraints.

\section{Is our treatment rigorous or a proxy?}
\label{sec:rigorous}

We must be precise here, because the catalogue tags it carefully. The
\emph{experimental} side of \code{L001} --- the MEG~II limit \eqref{eq:megii} ---
is fully \emph{rigorous}: a published, snapshotted, fact-checked number. The
\emph{prediction} side is an \emph{NDA estimate}, not a finite one-loop
computation: it uses the Perez--Randall dipole prefactor $4\times10^{-8}$ and a
fixed $\MKK$ convention rather than evaluating the loop integral over the actual
KK tower with the point's wavefunction overlaps. The adapter says so explicitly,
in the constant \code{LEPTON\_DIPOLE\_PROXY\_ASSUMPTION\_V1}, which is flagged
\code{NEEDS-HUMAN-PHYSICS}. So the correct tag is: \emph{rigorous experimental
limit, NDA/model prediction}. Where the lepton carrier is generated from a full
\code{YukawaResult} (not a caller-supplied stub) the catalogue marks the result
\code{rigorous} in the sense that no proxy input was injected --- but the
underlying dipole normalisation is still the NDA form, which is a model
assumption, not a derived loop.

These two statements are not in tension once one sees that they live on
\emph{orthogonal axes}, and the apparent paradox dissolves the moment they are
separated:
\begin{center}
\begin{tabular}{lll}
\toprule
 & \textbf{Axis 1: input provenance} & \textbf{Axis 2: physics fidelity} \\
 & (the \code{used\_proxy} flag) & (the dipole calculation) \\
\midrule
question asked & real \code{YukawaResult} or caller stub? & NDA estimate or finite loop? \\
``rigorous'' here & real $\Rightarrow$ no proxy input & --- (axis~1 only) \\
today's status    & can be \emph{rigorous} (real carrier) & always \emph{proxy} (NDA, never a loop) \\
\bottomrule
\end{tabular}
\end{center}
The catalogue's \code{rigorous} tag refers to \emph{axis~1 only}: it certifies
that the spurion came from a genuine \code{YukawaResult}, not a hand-supplied
stub. The \emph{physics} (axis~2) is \emph{proxy regardless} --- the dipole
normalisation is the NDA form for every point, real carrier or not. So a result
can be ``rigorous'' (axis~1) and ``proxy'' (axis~2) at once with no contradiction;
the gap to a truly rigorous prediction is entirely on axis~2 and is the subtlety
of \S\ref{part:subtle}.

% ====================================================================
\part{What is in OUR code}
\label{part:code}
% ====================================================================

\emph{Reader here for the physics may skim Part~\ref{part:code} on a first pass
and jump to the subtleties in Part~\ref{part:subtle}; this part documents the
exact code so the note doubles as an implementation reference.}

\section{The low-level dipole check: \code{flavorConstraints/muToEGamma.py}}
\label{sec:codemuegamma}

The physics primitive lives in \code{flavorConstraints/muToEGamma.py}. Its
module-level constants encode the Perez--Randall NDA convention exactly:
\begin{equation}
\code{PREFAC\_BR} = 4.0\times10^{-8},\quad
\code{BR\_LIMIT\_PAPER} = 1.2\times10^{-11},\quad
\code{C\_PAPER} = 0.02,\quad
\code{GAUGE\_KK\_ROOT\_NN} = 2.448687\ldots
\end{equation}
The two workhorse functions are:
\begin{itemize}[leftmargin=1.4em]
\item \code{coefficient\_from\_br\_limit(br\_limit, prefactor=PREFAC\_BR)} returns
$C=\sqrt{\BR_{\rm lim}/\text{prefactor}}$ --- the inversion of
\eqref{eq:spurbound}. With $\BR_{\rm lim}=1.2\times10^{-11}$ this gives
$\sqrt{1.2\times10^{-11}/4\times10^{-8}}\approx 0.0173$, which the paper rounds up
to $0.02$ (the source comment says exactly this); with the MEG~II
$1.5\times10^{-13}$ it gives $0.0019364916731037085$, the value stored as the
live \code{lfv\_C} in \code{L001.yaml}.
\item \code{check\_mu\_to\_e\_gamma\_raw(Y\_N\_bar, pmns, M\_KK, C=C\_PAPER,
reference\_scale=3000.0)} builds $\YNbmat=\Vpmns\cdot\mathrm{diag}(\YNb)$,
forms the Hermitian product $\YNbmat\YNbmat^\dagger$, extracts
the $(0,1)$ entry (the $e\mu$ spurion $\spur$), and compares
$\text{lhs}=|\spur|$ against $\text{rhs}=C\,(\MKK/3000)^2$. It returns
\code{lhs}, \code{rhs}, \code{passes} ($=\text{lhs}\le\text{rhs}$), \code{ratio},
\code{off\_diagonal\_12}, and the full \code{product\_matrix}.
\end{itemize}
Note the default \emph{internal} $\MKK$ convention: \code{check\_mu\_to\_e\_gamma}
falls back to \code{M\_KK = Lambda\_IR} (the geometric IR scale) unless an
explicit \code{M\_KK} is stored or overridden; the helper
\code{default\_m\_kk\_from\_lambda\_ir} exists to map $\LIR\mapsto x_1\LIR$ but is
``\emph{not} the repo's default LFV convention.'' This is the $2.45$ ambiguity
again (\S\ref{sec:Mconvention}).

\section{The LMFV carrier: \code{LMFVLeptonParameters}}
\label{sec:carrier}

The catalogue boundary is
\code{flavor\_catalog\_constraints/physics\_adapters/lepton.py}, which wraps the
primitive above. Its central object is the frozen dataclass
\code{LMFVLeptonParameters}, a read-only carrier for the Perez--Randall spurion.
It stores $\YN$, $\YNb$, the $3\times3$ matrices $\YNmat$ and
$\YNbmat$, the PMNS matrix, the precomputed \code{lmfv\_spurion}
$=\YNbmat\YNbmat^\dagger$, the scales $M_{KK}$, $M_N$, $v$,
$k$, $\LIR$, $\eps$, the bulk masses $c_L,c_E,c_N$, the ordering, the Majorana
phases, and \code{max\_abs\_ybar}. Its \code{\_\_post\_init\_\_} enforces the
physics as runtime invariants, which is the cleanest documentation of the model:
\begin{itemize}[leftmargin=1.4em]
\item the PMNS matrix must be \emph{unitary};
\item $\YNb = 2k\,\YN$ (the rescaling convention);
\item $\YNbmat = 2k\,\YNmat = \Vpmns\cdot\mathrm{diag}(\YNb)$
(the rotation \eqref{eq:YNrot});
\item \code{lmfv\_spurion} $= \YNbmat\YNbmat^\dagger$;
\item the \code{matching\_status} string may not contain any of
\code{\{"partial","deferred","proxy","recast"\}} --- a guard so that a carrier
built from a real \code{YukawaResult} cannot be silently tagged as a proxy.
\end{itemize}
The builder \code{lmfv\_lepton\_parameters\_from\_yukawa\_result} constructs this
carrier from a repo \code{YukawaResult} plus an explicit \code{m\_kk\_gev}, pulling
the PMNS from \code{neutrinos.neutrinoValues.get\_pmns(ordering, alpha, beta)}.
This is the ``locked'' path; the legacy \code{MuToEGammaProxyInput} (caller
supplies $\YNb$, PMNS, $\MKK$ directly) still works but is permanently flagged
\code{used\_proxy=True}.

\section{The catalogue constraint \code{L001}}
\label{sec:codeL001}

The catalogued constraint is
\code{flavor\_catalog\_constraints/primary/charged\_lepton/L001.py}. Its identity:
\begin{center}
\begin{tabular}{ll}
\toprule
\code{process\_id} & \code{"L001"} \\
\code{severity} & \code{Severity.HARD} \\
\code{observable} & \code{"BR(mu -> e gamma)"} \\
required extra & \code{"lepton\_lmfv\_parameters"} (an \code{LMFVLeptonParameters}) \\
optional extra & \code{"kk\_ew\_mass\_gev"} (consistency-checked, not used) \\
budget & MEG~II limit, $1.5\times10^{-13}$, from \code{L001.yaml} \\
reference scale & \code{3000.0}~GeV (Perez--Randall NDA convention) \\
\bottomrule
\end{tabular}
\end{center}
The numeric anchors are \emph{loaded from the YAML sidecar and
consistency-checked}, not hardcoded: \code{\_load\_l001\_anchor} asserts
that \code{repo\_default.br\_limit} equals the primary MEG~II limit and that
\code{repo\_default.lfv\_C} equals $\sqrt{\BR_{\rm lim}/\text{prefac}}$ to within
$10^{-12}$. The pure-new-physics branching fraction is
\begin{equation}
\BR_{\rm NP} = \code{prefactor\_br}\times|\spur|^2\times
   \left(\frac{3\ \mathrm{TeV}}{\MKK}\right)^4,
\label{eq:codeBR}
\end{equation}
computed by \code{mu\_to\_e\_gamma\_from\_lepton\_input} in the adapter (which
calls \code{check\_mu\_to\_e\_gamma\_raw} for the spurion and then squares it into
a branching fraction). The \code{ConstraintResult} reports
\code{predicted} $=\BR_{\rm NP}$, \code{sm\_prediction} $=0.0$,
\code{experimental} $=1.5\times10^{-13}$, \code{ratio} $=\BR_{\rm NP}/\BR_{\rm lim}$,
\code{budget} $=\BR_{\rm lim}$, and \code{passes} $=(\BR_{\rm NP}\le\BR_{\rm lim})$.

\section{The critical honesty point: the gate is BR, not $C$}
\label{sec:honesty}

This is the single most important implementation fact in the whole note, and it
is easy to get wrong because the literature (and our own docs) talk constantly
about ``the coefficient $C$.'' \textbf{In our code the veto is
$\BR_{\rm NP}\le\BR_{\rm lim}$. The coefficient $C=0.02$ does \emph{not} enter the
pass/fail decision at all.} It is computed, stored, and reported --- as
\code{lfv\_coefficient}, with the explicit diagnostic
\code{c\_lfv\_role = "dipole\_rhs\_diagnostic\_only"} --- but it only scales the
\emph{dipole right-hand side} \code{dipole\_rhs} $=C(\MKK/3000)^2$, which is a
diagnostic comparison, never the gate.

The proof is a unit test,
\code{test\_pass\_fail\_is\_pinned\_to\_br\_and\_independent\_of\_c\_lfv}. It runs
the \emph{same} benchmark carrier twice, once with \code{c\_lfv=0.02} (paper) and
once with \code{c\_lfv=100.0} (deliberately absurd), and asserts:
\begin{itemize}[leftmargin=1.4em,nosep]
\item the two \code{dipole\_rhs} values differ (so $C$ \emph{is} doing something
to the diagnostic);
\item the \emph{dipole} sub-check flips: \code{core\_passes} is \code{False} at
$C=0.02$ but \code{True} at $C=100$;
\item \emph{nevertheless} the two \code{branching\_fraction} values are
identical, the two \code{ratio\_to\_limit} values are identical, and
\code{paper.passes is loose.passes is False}.
\end{itemize}
In words: cranking the coefficient up by a factor of $5000$ flips the inert
dipole diagnostic from fail to pass, but the actual constraint verdict
\textbf{does not move}, because the verdict is the branching ratio against the
MEG~II limit. This was flagged by both reviewers in the W7 dual sign-off and is
recorded verbatim in the program ledger: \emph{``Veto $=$ BR\_NP $\le$ br\_limit
(MEG~II 1.5e-13); C$=$0.02 is an INERT diagnostic, NOT the gate.''} If you take
one thing from this note, take this.

\section{Severity, anchors, and the ``NOT EVALUATED'' path}
\label{sec:notinscan}

\code{L001} is \code{Severity.HARD}: a point that predicts $\BR_{\rm NP}>\BR_{\rm lim}$
is \emph{vetoed} outright. But there is a crucial subtlety about \emph{when} it is
evaluated. \code{L001.evaluate} first looks for the required extra
\code{"lepton\_lmfv\_parameters"} on the \code{ParameterPoint}. If it is absent
(or not an \code{LMFVLeptonParameters}), \code{L001} returns
\code{\_unevaluated\_result}: \code{passes=True}, \code{predicted=None}, with the
diagnostic \code{evaluated=False} and the note
``\emph{NOT EVALUATED --- no lepton dipole prediction available (LMFV lepton
carrier not on ParameterPoint)}'' and the semantics
``\emph{non-vetoing only; no BR(mu -> e gamma) prediction was evaluated}.''

This is exactly the state of the \emph{quark-only production scans}. The point
builder registers \code{"lepton\_lmfv\_parameters"} as a known extra
(\code{point\_builder.py}), but the quark-only builders never populate it ---
they build quark couplings, KK masses, and the EW spectrum, not a lepton carrier.
So although \code{L001} is live and \code{HARD}, in the 1M quark-only scans it is
a non-vetoing pass-through: \textbf{it is a lepton-sector constraint that is not
in the quark-only production scans}. To actually exercise it one builds a lepton
\code{ParameterPoint} carrying an \code{LMFVLeptonParameters} (e.g.\ via
\code{lmfv\_lepton\_parameters\_from\_yukawa\_result}), which is what the W7 tests
do and what a future lepton-sector scan would do.

\section{Provenance and status (W7, commit \code{8ca55c8})}
\label{sec:provenance}

The catalogue entry is heavily provenance-tracked. \code{L001.yaml} records the
MEG~II limit (snapshot
\code{megii2025\_arxiv2504\_15711.txt}, sha256 \code{3560d6\ldots}), the
paper-era reference (Perez--Randall 0805.4652, snapshot
\code{perez\_randall\_2008\_arxiv0805\_4652.txt}), and the repo defaults
(\code{prefac\_br=4.0e-8}, \code{c\_paper=0.02}, derived
\code{lfv\_C}$=0.0019364916731037085$). Its \code{status\_history} shows the full
dual-signoff trail: writer/checker cycles, an Opus arbitration
(\code{opus\_arbitration\_L001}) that ruled the $3$~TeV reference scale is a
theoretical convention rather than a measured observable, an Opus round-1
approval, and a fact-check verdict \code{VERIFIED} with zero findings. The W7
ledger entry summarises the wiring: the carrier
(``$\YN$/PMNS/$\MKK$ + spurion $\spur$'') wired to \code{L001} through the
existing \code{muToEGamma.py}, quark-only behaviour byte-identical (config hashes
\code{45e21a07585f7489}/\code{d96cb734f724aedb}), and the \code{L001.tex}/\code{.yaml}
sidecars left untouched.

% ====================================================================
\part{Subtleties and the gap to ``rigorous''}
\label{part:subtle}
% ====================================================================

\section{NDA, not a loop}
The prediction \eqref{eq:codeBR} is an NDA \emph{estimate} of the dipole, not a
finite one-loop calculation over the KK tower. As ABP emphasise, the dipole is
only genuinely calculable with a \emph{bulk} Higgs; with the brane-localised NDA
form the coefficient $4\times10^{-8}$ absorbs an $\mathcal{O}(1)$ uncertainty.
The adapter labels this honestly in
\code{LEPTON\_DIPOLE\_PROXY\_ASSUMPTION\_V1} (\code{NEEDS-HUMAN-PHYSICS}). A
rigorous upgrade would compute the dipole from the actual KK-fermion/KK-gauge
loop with the point's wavefunction overlaps, the same machinery that already
feeds the quark-sector $Z\bar b\bar b$ and $\Delta F=2$ constraints.

\section{The $C=0.02$ number and its rescaling}
\label{sec:Cnumber}
The coefficient $0.02$ is \emph{not} fundamental: it is
$\sqrt{\BR_{\rm lim}/4\times10^{-8}}$ evaluated at the \emph{2008} MEGA input
$1.2\times10^{-11}$ (giving $0.0173$, rounded up). At the MEG~II limit it becomes
$\approx 1.9\times10^{-3}$. The catalogue keeps \emph{both}: \code{c\_paper=0.02}
for paper reproduction, and the derived live \code{lfv\_C}$\approx 1.9\times10^{-3}$.
Because the bound on the spurion scales linearly with $C$, the MEG~II update
tightens the spurion bound by $\sqrt{1.2\times10^{-11}/1.5\times10^{-13}}\approx 9$
relative to the paper era --- equivalently the $\MKK$ floor climbs by $\sqrt{9}=3$.
The ledger records this as ``bounds rescale $\sim 3\times$ at MEG~II.''
\textbf{But remember \S\ref{sec:honesty}: in the live gate $C$ is inert; the
rescaling that actually matters is in the branching ratio, through $\BR_{\rm lim}$.}

\section{The $\MKK$ convention: the factor of $2.45$}
\label{sec:Mconvention}
``$\MKK$'' is overloaded, exactly as in the oblique-$S,T$ note. The physical
first KK gauge mass is $\MKK^{\rm phys}=x_1\LIR\approx 2.45\,\LIR$
(\code{GAUGE\_KK\_ROOT\_NN}). The repo's \emph{default LFV convention}, however,
uses $\MKK=\LIR$ (see the fallback in \code{check\_mu\_to\_e\_gamma} and the
docstring of \code{check\_mu\_to\_e\_gamma\_raw}). A bound quoted as ``$\sim10$~TeV''
in $\LIR$ units is ``$\sim 25$~TeV'' physical --- a $2.45\times$ difference, and
because the rate goes as $\MKK^{-4}$, a factor of $2.45^4\approx 36$ in the
branching ratio. Whichever convention the prefactor $4\times10^{-8}$ and the
$3$~TeV reference were calibrated in \emph{must} match the $\MKK$ fed to
\code{L001}. The helper \code{default\_m\_kk\_from\_lambda\_ir} exists precisely to
make the conversion explicit when one wants the physical convention.

\section{Dependence on the neutrino sector inputs}
The spurion $\spur$ is \emph{not} a free knob; it is fixed by the neutrino-sector
inputs through \eqref{eq:YNrot}: the PMNS angles and phases, the neutrino Yukawa
eigenvalues $Y_{N_i}$, the mass ordering, and (via the seesaw) the Majorana scale
$M_N$. Two consequences. First, the LFV bound and the neutrino-mass fit are
\emph{correlated}: one cannot tune the spectrum to fit oscillation data without
moving $\mu\to e\gamma$. Second, the Majorana phases (\code{majorana\_alpha},
\code{majorana\_beta}) and the ordering enter the carrier explicitly, so a scan
that marginalises over them is marginalising over the LFV prediction too. A
rigorous lepton scan must vary these consistently, not fix them.

\section{It is not (yet) in a production scan}
As \S\ref{sec:notinscan} explains, \code{L001} is live and \code{HARD} but
returns ``NOT EVALUATED'' wherever the lepton carrier is absent --- which is every
quark-only production scan to date. The constraint is wired and tested, but the
\emph{lepton-sector scan that would actually feed it} is follow-up scope (the
paper branch is quark-sector only). Until that scan exists, \code{L001}'s
practical effect on published results is null; its value today is as a tested,
ready-to-fire constraint and as documentation of the LMFV model choice.

\section{Checklist: what ``rigorous $\mu\to e\gamma$'' would require}
\begin{enumerate}[leftmargin=1.8em]
\item Replace the NDA prefactor \eqref{eq:codeBR} with a finite one-loop dipole
over the KK tower, using the point's wavefunction overlaps (bulk-Higgs scheme,
per ABP).
\item Fix a single, documented $\MKK$ convention (physical $x_1\LIR$ vs.\ $\LIR$)
and use it consistently in the prefactor calibration and the carrier.
\item Build a genuine lepton-sector scan that populates
\code{lepton\_lmfv\_parameters}, varying $\YN$, the PMNS angles/phases, the
ordering, and $M_N$ consistently with oscillation and cosmological data.
\item Decide whether to expose both a paper-reproduction mode ($C=0.02$, MEGA
input) and the live MEG~II mode, with the convention difference documented.
\item Cross-check the LMFV prediction against the anarchic and Chen--Yu
alternatives at the same benchmark, to bound the model-choice uncertainty.
\end{enumerate}

% ====================================================================
\part*{Annotated reading guide}
\addcontentsline{toc}{section}{Annotated reading guide}
% ====================================================================

\begin{description}[leftmargin=0em,style=nextline]
\item[Perez \& Randall, arXiv:0805.4652 (\emph{JHEP 01(2009)077}).] The model our
code adopts. Natural neutrino masses/mixings from warped geometry, the LMFV
hypothesis (bulk masses built from Yukawa spurions), and the NDA $\mu\to e\gamma$
dipole. Their Eq.~(29) (dipole operator) and Eqs.~(30)--(31) (NDA branching ratio
and spurion bound) are our \eqref{eq:prdipole}--\eqref{eq:prBR}, with $\YN\to\Vpmns
\,\mathrm{diag}(Y_N)$ their Eq.~(12); the coefficient $C\approx0.02$
\eqref{eq:spurbound} and the spurion $\spur$ are theirs.
Read \S\ref{sec:lmfv}--\S\ref{sec:brformula} alongside it. The single most
load-bearing reference for our \code{L001}.

\item[Agashe, Blechman, Petriello, hep-ph/0606021 (\emph{PRD 74 (2006) 053011}).]
The anarchic-RS LFV paper: LFV (especially $\mu\to e\gamma$) as a \emph{direct
probe} of the geometric origin of flavour, the RS-GIM mechanism in the lepton
sector, the bulk-vs-brane Higgs calculability issue, and the multi-TeV KK floor.
The reference for \S\ref{sec:mechanism}--\S\ref{sec:anarchic} and for why
custodial symmetry does not relax the LFV bound (\S\ref{sec:custodial}).

\item[MEG II Collaboration, arXiv:2504.15711 (2025).] The current primary
experimental limit \eqref{eq:megii}, $\BR<1.5\times10^{-13}$ (90\% C.L.). This is
the snapshotted, fact-checked anchor in \code{L001.yaml}. The rigorous side of the
constraint.

\item[MEG II 2024 (arXiv:2310.12614) and MEG 2016 (arXiv:1605.05081).] The
intermediate and prior limits ($7.5$, $3.1$, $4.2\times10^{-13}$) recorded in the
sidecar's experimental history; useful for understanding the rescaling of $C$
(\S\ref{sec:Cnumber}).

\item[Repo: \code{flavorConstraints/muToEGamma.py}.] The physics primitive:
\code{PREFAC\_BR}, \code{C\_PAPER}, \code{coefficient\_from\_br\_limit},
\code{check\_mu\_to\_e\_gamma\_raw}. Read \S\ref{sec:codemuegamma} with it open.

\item[Repo: \code{flavor\_catalog\_constraints/physics\_adapters/lepton.py}.] The
catalogue boundary and the \code{LMFVLeptonParameters} carrier (the spurion, the
invariants, the proxy flag). \S\ref{sec:carrier}.

\item[Repo: \code{.../primary/charged\_lepton/L001.py} and \code{L001.yaml}.] The
\code{Severity.HARD} constraint, the anchors, the ``NOT EVALUATED'' path, and the
\code{c\_lfv\_role = dipole\_rhs\_diagnostic\_only} diagnostic. \S\ref{sec:codeL001}.

\item[Repo: \code{tests/.../test\_L001.py},
\code{test\_pass\_fail\_is\_pinned\_to\_br\_and\_independent\_of\_c\_lfv}.] The
proof that the gate is the branching ratio, not $C$. The benchmark oracle
$\BR=1.55\times10^{-10}$, ratio $1033.9$. \S\ref{sec:honesty} and the worked
example.

\item[Ledger: \code{.orchestration/PHASE2\_PROGRAM\_LEDGER.md} (W7, commit
\code{8ca55c8}).] The wiring summary, the byte-identical-quark-scan guarantee, the
``$C$ inert'' note, and the LMFV-vs-alternatives verdict. \S\ref{sec:provenance}.
\end{description}

\bigskip
\noindent\rule{\linewidth}{0.4pt}\\
\footnotesize
\emph{Internal note. Equation \eqref{eq:prBR} is the NDA form our code uses; the
numerical inputs ($\code{PREFAC\_BR}=4\times10^{-8}$, $\code{C\_PAPER}=0.02$,
derived $\code{lfv\_C}\approx1.9\times10^{-3}$, MEG~II
$\BR_{\rm lim}=1.5\times10^{-13}$, reference scale $3$~TeV) are as stored in
\code{muToEGamma.py} and \code{L001.yaml}. The ``rigorous experimental limit /
NDA prediction'' tagging and the ``$C$ inert, BR is the gate'' fact follow the
code and the orchestration ledgers. Cross-check before citing externally.}

\end{document}
